模型上线之后，你迟早会被问一句``它到底学到了什么''。参数表在手边，几亿个数字一个也不会说话；换个随机种子重训，同样准确率的模型给出一组完全不同的数字。这个单元处理的就是这个处境：网络已经能算、能训，但关于它的一切陈述都还没有可以检验的形式。

前面几个单元讲的是机制——非线性带来什么、梯度怎样回传、优化器与归一化改变了什么。它们没有回答这些机制在具体任务上利用了哪一类结构。本单元也不承诺给出``深度网络为什么有效''的统一答案，它做的是把几种流行说法各自的适用范围划清楚，并给每一种配上可以推翻它的对照实验。

四篇各立一个判断。第一篇说明网络定义的是一个函数，参数只是它的一组坐标：置换单元编号、按 ReLU 的正齐次性缩放权重，都能在不改变函数的前提下改写全部参数，因此任何在这些变换下会变的陈述描述的是坐标而不是模型——``看参数''这条路从起点就是封的。第二篇处理核方法这条对照线：无限宽极限下梯度训练退化成固定核上的线性回归，中间层不再学习，于是 NTK 恰好给出``完全不学表示能到什么水平''的基线；深度模型区别于核方法的那部分，全在有限宽度里。第三篇检验流传最广的那条假设——数据位于低维流形附近——说明内在维数是尺度的函数而非数据的属性，以及在证据不足时该怎样把话说准。第四篇把解释本身形式化：敏感度、必要性与充分性是三个逻辑上互不蕴含的量，归因方法之间的不一致来自它们在回答不同的问题，而不是某一方算错了。

贯穿四篇的是同一条主线，也是整个 Part 的主线：深度模型把``选择表示''从人工设计变成可优化的对象，代价是几乎所有保证都退化为经验规律。在这样的对象上，理论类比和漂亮的可视化只能产生假设。可用的做法是先明确分析对象与所处的极限制度，再定义要测量的量，最后设计能推翻它的对照。

\subsection{怎么读这一章}

四篇分别从函数空间、核极限、低维结构与解释形式化切入，前两篇偏概念校准，后两篇偏检验流程。读的时候守住两条：不要把可视化图案直接当成机制解释；每一个关于模型的结论都要说清针对哪个输入、哪个输出、哪种扰动，以及是否经过稳定性与反事实检验。建议先熟悉 Part 13 的机器学习基础、Part 12 的流形与信息几何，以及 Part 14 前面几个单元。

\subsection{本章篇目}

以下篇目按概念依赖排列；若从具体问题进入，也可以先打开最接近当前困惑的一篇，再沿篇内链接回补。

\begin{itemize}
\tightlist
\item \hyperref[sec:14-Deep-Learning/10-Interpretability-Mathematics/01]{``神经网络的函数空间视角''}；
\item \hyperref[sec:14-Deep-Learning/10-Interpretability-Mathematics/02]{``特征学习与核方法的差异''}；
\item \hyperref[sec:14-Deep-Learning/10-Interpretability-Mathematics/03]{``高维数据的低维结构检验''}；
\item \hyperref[sec:14-Deep-Learning/10-Interpretability-Mathematics/04]{``模型解释的数学形式化''}。
\end{itemize}

读完这四篇不会让模型变得透明，但会让关于模型的陈述变得可以被反驳——在一个保证已经退化的领域里，可反驳是能守住的最高标准。
